\section{训练好了，模型怎样逐词生成文字？}
\label{intro:23-llm-systems}

\begin{idea}[title={同一个模型，在训练与生成时经历不同的计算流程}]
上一章已经完成了\textbf{因果掩码、错位标签、交叉熵与反向传播}。
现在暂时把可训练参数固定下来：如果输入是“春天、花”，
怎样让模型真正输出一个新的 Token？
\end{idea}

\subsection{预测分数还不是最终文字}

Transformer 的最后位置产生一条隐藏向量，
通过 \conceptlink{output-head}{词表输出头} 变成 $|\mathcal V|$ 个 \conceptlink{logits}{logits}。
这些分数并不是概率，必须经过 \conceptlink{softmax}{Softmax} 归一化。
比如只看三个候选，概率可能是
\[
p=(0.6,\ 0.3,\ 0.1).
\]
\textbf{贪心解码}总选择第一个候选；\textbf{采样}
则允许按这组概率随机选取一个。
选出的 Token ID 再由分词器的解码规则转换成文字。

\subsection{为什么有温度、Top-k 和 Top-p？}
\phantomsection\label{concept:sampling}

如果每次都只取最大概率的候选，
结果相对确定，却可能缺乏多样性。
一种做法是改变\textbf{候选的分布或选择范围}：
\begin{enumerate}
\item \textbf{温度（Temperature）}：把 logits 除以一个正数，
再做 Softmax。温度较低通常使分布更集中，温度较高则更平坦。
\item \textbf{Top-k}：保留概率最高的 $k$ 个候选，重新归一化再采样。
\item \textbf{Top-p}：从概率较大者开始累加，保留达到阈值的一组候选再采样。
\end{enumerate}
这些都是\textbf{推理时的选择策略}，
通常不需要重新训练网络权重，也不保证生成事实正确。

\subsection{为什么训练能一次算多个位置，生成却通常要逐个产生？}

训练时，“春天 / 花 / 开了 / 。”的正确 Token
已经在数据中，因此 \conceptlink{causal-mask}{因果掩码} 可以让各位置在同一次前向计算中
\textbf{同时预测它们各自的下一 Token}。

而生成时，模型先看“春天、花”，
在多个候选中选出下一 Token；
这个新 Token 成为下一步上下文的一部分，
才有条件选择之后的 Token：
\[
(t_1,\ldots,t_t)
\longrightarrow p(t_{t+1}\mid t_{\le t})
\longrightarrow t_{t+1}
\longrightarrow\text{继续预测}.
\]
这叫\textbf{自回归生成}。
它不是重复训练，而是反复使用已经学好的参数做前向计算。

\subsection{KV Cache：每生成一个 Token 都从头再算吗？}
\phantomsection\label{concept:kv-cache}

如果每生成一个 Token 都重新计算全部历史位置的 Key 和 Value，
许多运算会重复。因果 Self-Attention 中，
旧 Token 的 Key/Value 不需要随着新 Token 重新改变，
所以可以把它们存入\textbf{KV Cache} 供后续读取。

这样减少了重复计算，却引出新的问题：
\textbf{缓存越长，存储这些 K/V 需要的内存也越多。}
影响大小的量包括层数、序列长度、K/V 头数、
每头维度和数值精度。
因此，后面讲现代架构时再遇到 MQA、GQA，
你就能理解它们为什么关注\textbf{减少 K/V 头数}。

\subsection{上下文窗口不能当成永久记忆}

\textbf{上下文窗口}是一次模型计算允许使用的 Token 范围。
即使输入非常长，也不能保证模型无遗漏地利用每一处信息。
当历史超过限制时，应用可能截断、摘要或重新检索；
这些是应用层的策略，不能误认为模型参数在对话中被永久改写。

\begin{warningbox}
以上描述的是常见\textbf{自回归 Decoder-only 模型}。
不同系统可能使用推测解码、批处理或其他性能优化，
但这不改变“模型给出候选概率、系统选择 Token”的基本分工。
流畅的生成也不代表内容已经被外部证据核验。
\end{warningbox}

\begin{conclusion}
我们已经走完\textbf{表示、训练、Attention、Transformer、生成}
这一整条主线。新的疑问不再是“什么叫 Q/K/V”，
而是\textbf{如何降低生成成本、如何训练出更符合任务的行为，
以及怎样借助外部工具解决实际问题}。
这正是下一部分的起点。
\end{conclusion}
